% GENERATED FILE. Edit canonical lessons, then run npm run generate:books.
% canonical-source-hash: fnv1a64:de2b908b5a6a895a
% canonical-lessons: ML-C572-silikkuka, ML-C572-utayuka, ML-C572-vittujoli, ML-C572-atukkuka, ML-C572-marala, ML-R572-review-565-572, ML-R572-review-final-572

\chapter{Settling Into a Home}
\label{ch:ml-settling-into-a-home}
\hlchaptermodality{572}

\begin{chapteropening}
\textbf{By the end of this chapter:} \emph{I can say to get used to, to shatter, housework, to stack and a cobweb.}

Complete the last lesson of chapter 572: I can say to get used to, to shatter, housework, to stack and a cobweb.
\end{chapteropening}

\section[\texorpdfstring{\ml{ശീലിക്കുക} (śīlikkuka)}{ശീലിക്കുക (śīlikkuka)}]{\ml{ശീലിക്കുക} (śīlikkuka) — to get used to}

\label{lesson:ML-C572-silikkuka}



\begin{quote}
Before the new one: say the Malayalam for to hang up, then the Malayalam for to water plants.
\end{quote}

\subsection*{You'll want to know: \ml{ശീലിക്കുക}}
\textbf{\ml{ശീലിക്കുക}} — \emph{śīlikkuka} — \textquotedblleft{}to get used to\textquotedblright{}.

\begin{quote}
\ml{നേരത്തെ} \ml{എഴുന്നേൽക്കാൻ} \ml{ഞാൻ} \ml{ശീലിച്ചു}.
\end{quote}

\begin{quote}
\emph{nēratte eḻunnēlkkān ñān śīliccu.}
\end{quote}

\begin{quote}
I got used to getting up early.
\end{quote}

Past: \ml{ശീലിച്ചു}, \emph{śīliccu}.

From \ml{ശീലം}, a habit. What one gets used to doing takes -\ml{ആൻ}: \ml{എഴുന്നേൽക്കാൻ} \ml{ശീലിച്ചു}. In speech \ml{ശീലമായി}, it became a habit, is common.

\begin{cousinweb}
\textbf{\ml{ശീലിക്കുക}} (\emph{śīlikkuka}) is built from Sanskrit \emph{śīlayati}, \textquotedblleft{}practises repeatedly\textquotedblright{} (CDIAL 12502), the verb of \emph{śīla}, \textquotedblleft{}habit\textquotedblright{}, with the Malayalam verb ending -\ml{ിക്കുക}.
\end{cousinweb}

\begin{grammarlens}[title={What you've built}]
One more word for describing what you lived and felt.
\end{grammarlens}

\subsection*{Your turn}
\begin{itemize}
\raggedright
  \item \emph{Say it:} \emph{śīlikkuka}
  \item \emph{Say it:} \emph{śīlikkuka}, once more
  \item \emph{Say it:} \emph{nēratte eḻunnēlkkān ñān śīliccu.} — I got used to getting up early.
  \item \emph{Recall:} say the Malayalam for to hang up, then the Malayalam for to water plants, then say \emph{śīlikkuka} again
\end{itemize}

\begin{tcolorbox}[breakable,colback=teal!4,colframe=teal!35!black,title={Before you move on}]
What does \ml{ശീലിക്കുക} mean? (\textquotedblleft{}To get used to\textquotedblright{}.) Say it once more.
\end{tcolorbox}
\section[\texorpdfstring{\ml{ഉടയുക} (uṭayuka)}{ഉടയുക (uṭayuka)}]{\ml{ഉടയുക} (uṭayuka) — to shatter; to break into pieces}

\label{lesson:ML-C572-utayuka}



\begin{quote}
Before the new one: say the Malayalam for to water plants, then the Malayalam for to get used to.
\end{quote}

\subsection*{You'll want to know: \ml{ഉടയുക}}
\textbf{\ml{ഉടയുക}} — \emph{uṭayuka} — \textquotedblleft{}to shatter; to break into pieces\textquotedblright{}.

\begin{quote}
\ml{കയ്യിൽ} \ml{നിന്ന്} \ml{വീണ്} \ml{ഗ്ലാസ്} \ml{ഉടഞ്ഞു}.
\end{quote}

\begin{quote}
\emph{kayyil ninnŭ vīṇŭ glāsŭ uṭaññu.}
\end{quote}

\begin{quote}
The glass fell from my hand and shattered.
\end{quote}

Past: \ml{ഉടഞ്ഞു}, \emph{uṭaññu}.

Of glass, pots and eggs; \ml{ഉടയ്ക്കുക} is to break something. The past is \ml{ഉടഞ്ഞു}.

\begin{cousinweb}
\textbf{\ml{ഉടയുക}} (\emph{uṭayuka}) is native Dravidian: DEDR 946 gives the Malayalam verb as to break, as pots or nuts do, beside a transitive partner for breaking something to pieces (DEDR 946).
\end{cousinweb}

\begin{grammarlens}[title={What you've built}]
One more word for describing what you lived and felt.
\end{grammarlens}

\subsection*{Your turn}
\begin{itemize}
\raggedright
  \item \emph{Say it:} \emph{uṭayuka}
  \item \emph{Say it:} \emph{uṭayuka}, once more
  \item \emph{Say it:} \emph{kayyil ninnŭ vīṇŭ glāsŭ uṭaññu.} — The glass fell from my hand and shattered.
  \item \emph{Recall:} say the Malayalam for to water plants, then the Malayalam for to get used to, then say \emph{uṭayuka} again
\end{itemize}

\begin{tcolorbox}[breakable,colback=teal!4,colframe=teal!35!black,title={Before you move on}]
What does \ml{ഉടയുക} mean? (\textquotedblleft{}To shatter; to break into pieces\textquotedblright{}.) Say it once more.
\end{tcolorbox}
\section[\texorpdfstring{\ml{വീട്ടുജോലി} (vīṭṭujōli)}{വീട്ടുജോലി (vīṭṭujōli)}]{\ml{വീട്ടുജോലി} (vīṭṭujōli) — housework; chores}

\label{lesson:ML-C572-vittujoli}



\begin{quote}
Before the new one: say the Malayalam for to get used to, then the Malayalam for to shatter.
\end{quote}

\subsection*{You'll want to know: \ml{വീട്ടുജോലി}}
\textbf{\ml{വീട്ടുജോലി}} — \emph{vīṭṭujōli} — \textquotedblleft{}housework; chores\textquotedblright{}.

\begin{quote}
\ml{ഞായറാഴ്ച} \ml{എല്ലാവരും} \ml{വീട്ടുജോലിയിൽ} \ml{സഹായിച്ചു}.
\end{quote}

\begin{quote}
\emph{ñāyaṟāḻca ellāvaruṁ vīṭṭujōliyil sahāyiccu.}
\end{quote}

\begin{quote}
On Sunday everyone helped with the housework.
\end{quote}

\ml{വീട്} house + \ml{ജോലി} work; \ml{വീട്} becomes \ml{വീട്ടു}- in a compound.

\begin{cousinweb}
\textbf{\ml{വീട്ടുജോലി}} (\emph{vīṭṭujōli}) is a Malayalam compound built from \ml{വീട്}, \textquotedblleft{}a house\textquotedblright{}, which is native (DEDR 5393), and \ml{ജോലി}, \textquotedblleft{}work\textquotedblright{}.
\end{cousinweb}

\begin{grammarlens}[title={What you've built}]
One more word for describing what you lived and felt.
\end{grammarlens}

\subsection*{Your turn}
\begin{itemize}
\raggedright
  \item \emph{Say it:} \emph{vīṭṭujōli}
  \item \emph{Say it:} \emph{vīṭṭujōli}, once more
  \item \emph{Say it:} \emph{ñāyaṟāḻca ellāvaruṁ vīṭṭujōliyil sahāyiccu.} — On Sunday everyone helped with the housework.
  \item \emph{Recall:} say the Malayalam for to get used to, then the Malayalam for to shatter, then say \emph{vīṭṭujōli} again
\end{itemize}

\begin{tcolorbox}[breakable,colback=teal!4,colframe=teal!35!black,title={Before you move on}]
What does \ml{വീട്ടുജോലി} mean? (\textquotedblleft{}Housework; chores\textquotedblright{}.) Say it once more.
\end{tcolorbox}
\section[\texorpdfstring{\ml{അടുക്കുക} (aṭukkuka)}{അടുക്കുക (aṭukkuka)}]{\ml{അടുക്കുക} (aṭukkuka) — to stack; to arrange}

\label{lesson:ML-C572-atukkuka}



\begin{quote}
Before the new one: say the Malayalam for to shatter, then the Malayalam for housework.
\end{quote}

\subsection*{You'll want to know: \ml{അടുക്കുക}}
\textbf{\ml{അടുക്കുക}} — \emph{aṭukkuka} — \textquotedblleft{}to stack; to arrange\textquotedblright{}.

\begin{quote}
\ml{അവൾ} \ml{പുസ്തകങ്ങൾ} \ml{ഷെൽഫിൽ} \ml{അടുക്കി}.
\end{quote}

\begin{quote}
\emph{avaḷ pustakaṅṅaḷ ṣelphil aṭukki.}
\end{quote}

\begin{quote}
She stacked the books on the shelf.
\end{quote}

Past: \ml{അടുക്കി}, \emph{aṭukki}.

\ml{അടുക്ക്}, a pile, a layer, is the noun.

\begin{cousinweb}
\textbf{\ml{അടുക്കുക}} (\emph{aṭukkuka}) is native Dravidian: DEDR 80 lists \ml{അടുക്ക്}, \textquotedblleft{}a row, a pile, a layer\textquotedblright{}, and the verb, to pile up (DEDR 80).
\end{cousinweb}

\begin{grammarlens}[title={What you've built}]
One more word for describing what you lived and felt.
\end{grammarlens}

\subsection*{Your turn}
\begin{itemize}
\raggedright
  \item \emph{Say it:} \emph{aṭukkuka}
  \item \emph{Say it:} \emph{aṭukkuka}, once more
  \item \emph{Say it:} \emph{avaḷ pustakaṅṅaḷ ṣelphil aṭukki.} — She stacked the books on the shelf.
  \item \emph{Recall:} say the Malayalam for to shatter, then the Malayalam for housework, then say \emph{aṭukkuka} again
\end{itemize}

\begin{tcolorbox}[breakable,colback=teal!4,colframe=teal!35!black,title={Before you move on}]
What does \ml{അടുക്കുക} mean? (\textquotedblleft{}To stack; to arrange\textquotedblright{}.) Say it once more.
\end{tcolorbox}
\section[\texorpdfstring{\ml{മാറാല} (māṟāla)}{മാറാല (māṟāla)}]{\ml{മാറാല} (māṟāla) — a cobweb}

\label{lesson:ML-C572-marala}



\begin{quote}
Before the new one: say the Malayalam for housework, then the Malayalam for to stack.
\end{quote}

\subsection*{You'll want to know: \ml{മാറാല}}
\textbf{\ml{മാറാല}} — \emph{māṟāla} — \textquotedblleft{}a cobweb\textquotedblright{}.

\begin{quote}
\ml{പഴയ} \ml{മുറിയുടെ} \ml{മൂലകളിൽ} \ml{മാറാല} \ml{ഉണ്ടായിരുന്നു}.
\end{quote}

\begin{quote}
\emph{paḻaya muṟiyuṭe mūlakaḷil māṟāla uṇṭāyirunnu.}
\end{quote}

\begin{quote}
There were cobwebs in the corners of the old room.
\end{quote}

Old cobwebs full of dust; a fresh web is \ml{ചിലന്തിവല}.

\begin{cousinweb}
\textbf{\ml{മാറാല}} (\emph{māṟāla}) is a household word; no source we can open explains it, so its origin is not settled.
\end{cousinweb}

\begin{grammarlens}[title={What you've built}]
One more word for describing what you lived and felt.
\end{grammarlens}

\subsection*{Your turn}
\begin{itemize}
\raggedright
  \item \emph{Say it:} \emph{māṟāla}
  \item \emph{Say it:} \emph{māṟāla}, once more
  \item \emph{Say it:} \emph{paḻaya muṟiyuṭe mūlakaḷil māṟāla uṇṭāyirunnu.} — There were cobwebs in the corners of the old room.
  \item \emph{Recall:} say the Malayalam for housework, then the Malayalam for to stack, then say \emph{māṟāla} again
\end{itemize}

\begin{tcolorbox}[breakable,colback=teal!4,colframe=teal!35!black,title={Before you move on}]
What does \ml{മാറാല} mean? (\textquotedblleft{}A cobweb\textquotedblright{}.) Say it once more.
\end{tcolorbox}
\section[(dialogue)]{Words for what you lived and felt — the last review}

\label{lesson:ML-R572-review-565-572}



\begin{quote}
Say the Malayalam for a painter and a musician.
\end{quote}

\subsection*{Your turn}
Say these aloud:

\begin{itemize}
\raggedright
  \item a painter, a musician, a dancer, a writer, to produce a film
  \item clay, chess, a card game, gardening, to meditate
  \item to go through, to confront, to hold out, to overcome, to adapt
  \item an athlete, a coach, a goal, a referee, to compete
  \item a sports champion, an opponent, a tournament, to advance, to fall behind
  \item furniture, a balcony, a tenant, to do the laundry, to polish
  \item repairs, a dustbin, a living room, to hang up, to water plants
  \item to get used to, to shatter, housework, to stack, a cobweb
\end{itemize}

\begin{tcolorbox}[breakable,colback=teal!4,colframe=teal!35!black,title={Before you move on}]
A painter? (\textbf{\ml{ചിത്രകാരൻ}}, \emph{citrakāran}.) A musician? (\textbf{\ml{സംഗീതജ്ഞൻ}}, \emph{saṁgītajñan}.) A dancer? (\textbf{\ml{നർത്തകി}}, \emph{narttaki}.) A writer? (\textbf{\ml{എഴുത്തുകാരൻ}}, \emph{eḻuttukāran}.) To produce a film? (\textbf{\ml{നിർമ്മിക്കുക}}, \emph{nirmmikkuka}.) Clay? (\textbf{\ml{കളിമണ്ണ്}}, \emph{kaḷimaṇṇŭ}.) Chess? (\textbf{\ml{ചെസ്സ്}}, \emph{cessŭ}.) A card game? (\textbf{\ml{ചീട്ടുകളി}}, \emph{cīṭṭukaḷi}.) Gardening? (\textbf{\ml{തോട്ടപ്പണി}}, \emph{tōṭṭappaṇi}.) To meditate? (\textbf{\ml{ധ്യാനിക്കുക}}, \emph{dhyānikkuka}.) To go through? (\textbf{\ml{അനുഭവിക്കുക}}, \emph{anubhavikkuka}.) To confront? (\textbf{\ml{നേരിടുക}}, \emph{nēriṭuka}.) To hold out? (\textbf{\ml{പിടിച്ചുനിൽക്കുക}}, \emph{piṭiccunilkkuka}.) To overcome? (\textbf{\ml{മറികടക്കുക}}, \emph{maṟikaṭakkuka}.) To adapt? (\textbf{\ml{പൊരുത്തപ്പെടുക}}, \emph{poruttappeṭuka}.) An athlete? (\textbf{\ml{കായികതാരം}}, \emph{kāyikatāraṁ}.) A coach? (\textbf{\ml{പരിശീലകൻ}}, \emph{pariśīlakan}.) A goal? (\textbf{\ml{ഗോൾ}}, \emph{gōḷ}.) A referee? (\textbf{\ml{റഫറി}}, \emph{ṟaphaṟi}.) To compete? (\textbf{\ml{മത്സരിക്കുക}}, \emph{matsarikkuka}.) A sports champion? (\textbf{\ml{ചാമ്പ്യൻ}}, \emph{cāmpyan}.) An opponent? (\textbf{\ml{എതിരാളി}}, \emph{etirāḷi}.) A tournament? (\textbf{\ml{ടൂർണമെന്റ്}}, \emph{ṭūrṇamenṟŭ}.) To advance? (\textbf{\ml{മുന്നേറുക}}, \emph{munnēṟuka}.) To fall behind? (\textbf{\ml{പിന്നിലാകുക}}, \emph{pinnilākuka}.) Furniture? (\textbf{\ml{ഫർണിച്ചർ}}, \emph{pharṇiccar}.) A balcony? (\textbf{\ml{ബാൽക്കണി}}, \emph{bālkkaṇi}.) A tenant? (\textbf{\ml{വാടകക്കാരൻ}}, \emph{vāṭakakkāran}.) To do the laundry? (\textbf{\ml{അലക്കുക}}, \emph{alakkuka}.) To polish? (\textbf{\ml{മിനുക്കുക}}, \emph{minukkuka}.) Repairs? (\textbf{\ml{അറ്റകുറ്റപ്പണി}}, \emph{aṟṟakuṟṟappaṇi}.) A dustbin? (\textbf{\ml{ചവറ്റുകുട്ട}}, \emph{cavaṟṟukuṭṭa}.) A living room? (\textbf{\ml{സ്വീകരണമുറി}}, \emph{svīkaraṇamuṟi}.) To hang up? (\textbf{\ml{തൂക്കിയിടുക}}, \emph{tūkkiyiṭuka}.) To water plants? (\textbf{\ml{നനയ്ക്കുക}}, \emph{nanaykkuka}.) To get used to? (\textbf{\ml{ശീലിക്കുക}}, \emph{śīlikkuka}.) To shatter? (\textbf{\ml{ഉടയുക}}, \emph{uṭayuka}.) Housework? (\textbf{\ml{വീട്ടുജോലി}}, \emph{vīṭṭujōli}.) To stack? (\textbf{\ml{അടുക്കുക}}, \emph{aṭukkuka}.) A cobweb? (\textbf{\ml{മാറാല}}, \emph{māṟāla}.)
\end{tcolorbox}
\section[(dialogue)]{Settling Into a Home — once more}

\label{lesson:ML-R572-review-final-572}



\begin{quote}
Say the Malayalam for to get used to and to shatter.
\end{quote}

\subsection*{Your turn}
\begin{itemize}
\raggedright
  \item \emph{Say it:} to get used to, to shatter, housework, to stack, a cobweb
\end{itemize}

\begin{tcolorbox}[breakable,colback=teal!4,colframe=teal!35!black,title={Before you move on}]
To get used to? (\textbf{\ml{ശീലിക്കുക}}, \emph{śīlikkuka}.) To shatter? (\textbf{\ml{ഉടയുക}}, \emph{uṭayuka}.) Housework? (\textbf{\ml{വീട്ടുജോലി}}, \emph{vīṭṭujōli}.) To stack? (\textbf{\ml{അടുക്കുക}}, \emph{aṭukkuka}.) A cobweb? (\textbf{\ml{മാറാല}}, \emph{māṟāla}.)
\end{tcolorbox}
